\documentclass[11pt]{article}
\usepackage[utf8]{inputenc}
\usepackage[T1]{fontenc}
\usepackage{lmodern}
\usepackage{textcomp}
\usepackage{amsmath, amssymb, amsthm}
\usepackage{graphicx}
\usepackage{listings}
\usepackage[hidelinks]{hyperref}

\theoremstyle{plain}
\newtheorem{theorem}{Teorema}[section]
\newtheorem{lemma}[theorem]{Lema}
\newtheorem{proposition}[theorem]{Proposición}
\newtheorem{corollary}[theorem]{Corolario}

\theoremstyle{definition}
\newtheorem{definition}[theorem]{Definición}


\title{c3}
\date{}
\begin{document}
\maketitle

\section{Chapter 3}

\section{Essential Dictionary II}

We further develop our mathematical dictionary, introducing the words sequence, sum, and equation. This will enable us to describe mathematical expressions of respectable complexity. As in the previous chapter, the last section contains advanced terms which are not essential for the rest of this book.

\section{3.1 Sequences}

A sequence is an ordered list of objects, not necessarily distinct, called the terms (or the elements) of the sequence. The terms of a sequence are represented by a common symbol, and each term is identified by an integer subscript:

(a1, a2, . . . , an) (a1, a2, . . . ). (3.1)

Here the common symbol is a, and the integer values assumed by the subscript begin from 1. The quantity a1 reads ``a sub 1'', etc. Subscripts may also begin from 0, or from anywhere. The expression on the left denotes a finite sequence, the one on the right suggests that the sequence is infinite.

The length of a sequence is the number of its elements. Two sequences are equal if they have the same length, and if the corresponding terms are equal. If k is an unspecified integer, then ak is called the general term of the sequence.

For example, the sequence of primes

(p1, p2, p3, . . . ) = (2, 3, 5, . . .)

is infinite. The general term pk is the kth prime number.

There are other notations for sequences, displaying the general term alongside information about the subscript range:

© Springer-Verlag London 2014 F. Vivaldi, Mathematical Writing, Springer Undergraduate Mathematics Series, DOI 10.1007/978-1-4471-6527-9\_3

31

32 3 Essential Dictionary II

(ak)$\infty$

(ak)n

(ak)n

k=1 (ak)k$\geq$1 (ak). (3.2)

k=1

1

There are also the doubly-infinite sequences, where the subscript runs through all the integers:

(ak)$\infty$

k=$-$$\infty$= (. . . , a$-$1, a0, a1, . . .).

In Sect. 6.2 we shall discuss the usage of these notations.

The ellipsis is ubiquitous in sequence notation. When we use it, we must make sure that the missing terms are defined unambiguously. Thus the general term of the sequence of monomials

(2x, 2x2, 2x3, . . .)

is clearly equal to 2xk. However, the expression

(3, 5, 7, . . .)

is ambiguous, because there are several plausible alternatives for the identity of the omitted terms, such as (9, 11, 13, . . .) or (11, 13, 17, . . .). In the former case, we resolve the ambiguity by displaying the general term:

(3, 5, . . . , 2k + 1, . . .).

In the latter, we need an accompanying sentence.

A sub-sequence of a sequence (ak) is any sequence obtained from (ak) by deleting terms. For instance, the primes that give remainder 1 upon division by 4 form a sub-sequence of the sequence of primes:

(5, 13, 17, 29, . . .).

Some types of finite sequences have a dedicated terminology. A two-element sequence is an (ordered) pair, and a three-element sequence a triple. Occasionally one sees the terms quadruple or quintuple (I wouldn't go much beyond that), while an n-element sequence may be called an n-uple. A finite sequence of numbers may be called a vector, in which case we speak of dimension rather than length.

An infinite sequence (a1, a2, . . . ) represents a function defined over the natural numbers. If the elements of the sequence belong to a set A, then such a function is defined as

a : N $\to$A k $\to$ak.

We see that in the expression ak, the symbol a is the function's name, the subscript k is an element of the domain, and ak $\in$A is the value a(k) of the function at k. This

3.1 Sequences 33

interpretation clarifies the meaning of expressions such as ak2: it is the composition of two functions, much like sin(x2).

Many familiar constructs involve sequences. Let A be a set of numbers, and let (a0, . . . , an) be a finite sequence of elements of A with an = 0. A polynomial over

A in the indeterminate x is an expression of the type

a0 + a1x + a2x2 + \textperiodcentered{} \textperiodcentered{} \textperiodcentered{} + anxn. (3.3)

In this notation, repeated addition is represented by the raised ellipsis. The elements of the sequence are called the coefficients of the polynomial, and the integer n is its degree. The coefficients a0 and an are called, respectively, the constant and the leading coefficient. Each addendum in a polynomial is called a monomial, and a polynomial with two terms is a binomial. A polynomial of degree two is said to be quadratic; then we have cubic, quartic, and quintic polynomials. The set of all polynomials over the set A with indeterminate x is denoted by A[x]. For example

x2 $-$x $-$1 $\in$Z[x] 2 $-$y3 $\in$Q[y]. 1

A rational function is the ratio of two polynomials:

a0 + a1x + \textperiodcentered{} \textperiodcentered{} \textperiodcentered{} + anxn

b0 + b1x + \textperiodcentered{} \textperiodcentered{} \textperiodcentered{} + bmxm . (3.4)

Its degree is the largest of m and n (assuming that anbm = 0). The set of all rational functions with coefficients in a set A and indeterminate x is denoted by A(x).

A multivariate polynomial is a polynomial in more than one indeterminate. (The term univariate is used to differentiate from multivariate.)

x2y2 $-$1

2 x4 $-$xy3 $\in$Q[x, y].

The total degree of each monomial is the sum of the degrees of the indeterminates, and the degree of a polynomial is the largest total degree among the monomials with non-zero coefficient. A multivariate polynomial is homogeneous if all monomials have the same total degree. The expression above may be described as

A homogeneous quartic polynomial in two indeterminates with rational coefficients.

Example. Explain what is a polynomial. [ $\varepsilon$ ]

A polynomial is a finite sum. Each term, called a monomial, is the product of a coefficient (typically, a real or complex number) and one or more indeterminates, each raised to some non-negative integer power.

34 3 Essential Dictionary II

\section{3.2 Sums}

Sums can be found in every corner of mathematics. We develop the dictionary and the very rich notation associated to summation.

Given a finite sequence of numbers (a1, . . . , an), we form the sum of its elements:

n

ak = a1 + a2 + \textperiodcentered{} \textperiodcentered{} \textperiodcentered{} + an. (3.5)

k=1

This notation was introduced by Fourier.1 It is called the (delimited) sigma-notation, as it makes use of the capital Greek letter with that name, called the summation symbol. The subscript k is the index of summation, while 1 and n are, respectively, the lower limit and upper limit of summation. The quantity ak is the general term of the sum. The integer sequence (1, 2, . . . , n), specifying the values assumed by the index of summation, is called the range of summation.

The summation index is a dummy variable. This a variable used for internal book-keeping (like an integration variable), and its identity does not affect the value of the sum:

n

n

k2 =

j2 = 12 + 22 + \textperiodcentered{} \textperiodcentered{} \textperiodcentered{} + n2.

k=1

j=1

I suggest that you adopt a summation symbol (one of i, j, k,l, m, n, see Sect. 6.2) and stick to it, unless there is a good reason to change it.

There are variants to the delimited sigma-notation (3.5). In an unrestricted sum, range information may be omitted altogether:

n

= 2n n $\geq$0.

k

k

(This sum has in fact only finitely many non-zero terms.) The summation range may also be specified by inequalities placed below the summation symbols:

a j,k. (3.6)

ak

ak

1$\leq$k$\leq$n

k$\geq$1

1$\leq$j,k$\leq$N

Further conditions may be added to alter the range of summation:

ak = a$-$2 + a$-$1 + a1 + a2

0<|k|$\leq$2

1 Jean Baptiste Joseph Fourier (French: 1768--1830).

3.2 Sums 35

ak = a1 + a5 + a7 + a11

1$\leq$k$\leq$12 gcd(k,12)=1

ak = a2 + a3 + a5 + a7.

1$\leq$k$\leq$10 k prime

The advantage of this notation---called the standard form of the sigma-notation---is that the summation index is no longer restricted to a sequence of consecutive integers.

Example. Consider the following manipulation:

2k+2 =

2k+2 =

2 j = 2n $-$1. (3.7)

$-$2$\leq$k$\leq$n$-$3

0$\leq$k+2$\leq$n$-$1

0$\leq$j$\leq$n$-$1

After adding 2 to each term in the inequalities, we have simply replaced k + 2 with

j, obtaining the sum of a geometric progression which is evaluated explicitly. With this notation the change of summation index is unproblematic.

Example. For any natural number n, let $\phi$(n) be the number of positive integers smaller than n and relatively prime to it (with $\phi$(1) = 1). Thus $\phi$(12) = 4. This is Euler's $\phi$-function2 of number theory, which is defined in symbols as follows:

1, n > 1. (3.8)

$\phi$(1) = 1 $\phi$(n) =

1$\leq$k<n gcd(k,n)=1

Simply by letting ak = 1 in (3.5), we have turned a summation into an algorithm that counts the elements of the set specified by the given conditions. This neat device illustrates the flexibility of the sigma-notation. Alternatively, Euler's function may be defined as the cardinality of a set, using the cardinality symbol `\#' (see Sect. 2.1):

$\phi$(1) = 1 $\phi$(n) = \#\{k $\in$N : k < n, gcd(k, n) = 1\} n > 1.

Sums may be nested. A double sum is defined as follows:

K

J

J

K

def =

a j,k

a j,k

j=1

k=1

j=1

k=1

K

K

K

=

a1,k +

a2,k + \textperiodcentered{} \textperiodcentered{} \textperiodcentered{} +

aJ,k.

k=1

k=1

k=1

2 Leonhard Euler (Swiss: 1707--1783). An accessible account of Euler's mathematics is found in [11].

36 3 Essential Dictionary II

The sum in parentheses is a function of the outer summation index j; this sum is performed repeatedly, each time with a different value of j. The use of matching symbols for the index and the upper limit of summation ( j, J, k, K) is particularly appropriate here. If the two ranges of summations are independent, inner and outer sums can be swapped. The commutative and associative laws of addition ensure that the value of the sum will not change. We illustrate this process with an example.

1

3

a j,k = (a0,1 + a0,2 + a0,3) + (a1,1 + a1,2 + a1,3)

j=0

k=1

= (a0,1 + a1,1) + (a0,2 + a1,2) + (a0,3 + a1,3)

3

1

a j,k.

=

k=1

j=0

This chain of equalities adopts a standard layout and alignment. If the indices in a double sum have the same range, then they may be grouped together:

N

N

N

ai, j =

ai, j. (3.9)

i, j=1

i=1

j=1

A sum with infinitely many summands is called a series:

$\infty$

ak = a1 + a2 + \textperiodcentered{} \textperiodcentered{} \textperiodcentered{} . (3.10)

k=1

The summation range may be specified in several ways:

$\infty$

ak.

ak

ak

ak

k$\geq$1

k=$-$$\infty$

k$\in$Z

k

We haven't yet explained what (3.10) means. We write the definition using the assignment operator:

$\infty$

n

$\nabla$= lim

ak. (3.11)

ak

n$\to$$\infty$

k=1

k=1

The above limit---if it exists---is called the sum of the series, and the series is said to converge. Otherwise the series diverges. If a series has non-negative terms, then convergence is sometimes expressed with the suggestive notation (cf. (2.3), p. 11)

ak < $\infty$. (3.12)

k$\geq$0

3.2 Sums 37

Apowerseriesisaserieswhosetermsfeatureincreasingpowersofanindeterminate:

akxk. (3.13)

k$\geq$0

For the values of x for which the series (3.13) converges, the power series represents a function.

The terminology and notation introduced for sums extends with obvious modifications to products:

n

ak = a1 \textperiodcentered{} a2 \textperiodcentered{} \textperiodcentered{} \textperiodcentered{} an. (3.14)

k=1

An infinite product is called just that---there is no special name for it. Its value is defined as the limit of a sequence of finite products, as we did for sums (3.11). Infinite products are often written in the form

(1 + ak)

k$\geq$0

because for convergence we must have ak $\to$0.

All constructs introduced above may be extended to combinations of sums and products.

Exercise 3.1 Write out the following sums in full.

ak 2.

1.

a$-$k 3.

a1$-$k

0$\leq$k$-$1<3

k2<9

k2$\leq$k+2

ak 6.

4.

ak 5.

ak2

|k$-$3|<5 gcd(k,6)>1

k$\in$2Z+1

k2$\leq$9

|k|<5

\section{3.3 Equations and Identities}

Let f and g be functions with the same domain X and co-domain Y. An equation (on or over X) is an expression of the type

f (x) = g(x). (3.15)

The quantity x is the equation's unknown.

x2 $-$2x $-$4 = 0 An algebraic equation. cos(x) = sin(x) A trigonometric equation. log(1 + x) = $-$x A transcendental equation.

(3.16)

38 3 Essential Dictionary II

These are equations over (a subset of) R, and each equation is described in broad terms by an attribute (algebraic, trigonometric, transcendental). We shall learn more about these attributes in Sect. 3.4.2.

The expression (3.15) defines a property that each point x $\in$X either has or doesn't have.3 This prompts the definition of the solution set of Eq. (3.15), given by

\{x $\in$X : f (x) = g(x)\}.

For instance, over the real numbers the solution set of the first equation in (3.16) is \{1 +

$\surd$ 5, 1 $-$

$\surd$

5\}, while over Q the solution set is empty. We see that the solution set of an equation depends on the ambient set.

If the co-domain Y of f and g is a set of numbers (e.g., Y = Z, Q, R, C), then, by replacing f by f $-$g, we can reduce Eq. (3.15) to the simpler form

f (x) = 0. (3.17)

An element x of the solution set of this equation is called a zero of f , but if f (x) is a polynomial, we speak of a root of f . We also say that f vanishes at x. A function f vanishes identically on a set if it vanishes at every point of this set, in which case we use the emphatic notation f (x) $\equiv$0. For example, the real function x $\to$sin($\pi$x) vanishes identically on Z.

More generally, an equation may be reduced to the form (3.17) if the co-domain of f and g is a group with respect to addition. In this case the zero on the right-hand side is the zero element of the group. This is not necessarily the number 0, but it could be a zero matrix, a zero polynomial, a zero function, etc.

An example is given by the differential equations, which relate a function to its derivatives:

dt $-$x = 0 t2 dx2

dt2 + t dx

dx

dt + (t2 $-$1)x = 0. (3.18)

Here the unknown is x (not t), where x = x(t) is a function, and the underlying ambient set is a set of differentiable functions. The symbol 0 represents the zero function, namely the constant function that assumes the value zero everywhere. We'll consider other types of equations in Sect. 3.5.

An identity (or indeterminate equation) is an equation whose solution set is equal to the ambient set:

(x $-$1)3 = x3 $-$3x2 + 3x $-$1.

After simplification, every identity reduces to the standard form 0 = 0. This doesn't mean that identities are trivial, far from it; identities express equivalence of functions. However, they are ephemeral quantities, which disappear if they are simplified.

3 In Chap. 4 we shall see that an equation is a special type of predicate, which in turn is a special type of function.

3.3 Equations and Identities 39

Example. The identity

x2k + y2k

n$-$1

x2n $-$y2n = (x $-$y)

k=0

gives the full factorisation of the difference of two monomials whose degree is a power of 2, into the product of polynomials with integer coefficients.

Example. Over the set R2, the expression x + y = y +x is an identity, representing the commutativity of the addition of real numbers. The similar expression x + y = 1 $-$y is an equation, whose solution set is a line in R2.

By restricting the ambient set to the solution set, every equation becomes an identity. Clearly, sin($\pi$x) = 0 is an equation over R and an identity over Z. For a more subtle example, consider the equation f (x) = x5 $-$x = 0. The factorisation x5$-$x = x(x$-$1)(x+1)(x2+1) shows that the solution set over C is \{0, $\pm$1, $\pm$$\surd$$-$1\},

a set with five elements. Consider now the set

X = Z/5Z = \{0 + 5Z, 1 + 5Z, 2 + 5Z, 3 + 5Z, 4 + 5Z\}

of congruence classes modulo 5. Let us evaluate our function f at all points of Z/5Z, writing k for k + 5Z:

f (0) = 0 $\equiv$0 (mod 5) f (1) = 0 $\equiv$0 (mod 5)

f (2) = 30 $\equiv$0 (mod 5) f (3) = 240 $\equiv$0 (mod 5)

f (4) = 1020 $\equiv$0 (mod 5).

The function f vanishes identically over Z/5Z, and hence the expression x5 = x is an identity!

An expression of the type

 

f1(x) = 0 f2(x) = 0

x = (x1, . . . , xm) (3.19)

... fn(x) = 0



where all functions have the same domain and co-domain, is called a system of n simultaneous equations in m unknowns. The solution set of a system of equations is the intersection of the solution sets of the individual equations.

40 3 Essential Dictionary II

Example. Explain what is an equation, and its solutions. [ $\varepsilon$ ]

Bad: An equation is when we equate two functions. The solution is when the func-

tions are the same.

The inappropriate use of `when' is easily spotted (see Sect. 1.1), but there is a more serious flaw. The expression `equating two functions' means that we seek conditions under which the two functions become the same function. That's not what we had in mind. The operands of the equal sign in expression (3.15) are not functions, but rather values of functions.

Good: An equation is an expression that identifies the value of two functions at a

generic point of their common domain. The solutions of an equation are the points at which the two functions assume the same value.

(The expression `equating two functions' may be appropriate for functional equations, see (3.21).)

\section{3.4 Expressions}

The generic term expression indicates the symbolic encoding of a mathematical object. For instance, the string of symbols `2 + 3' is a valid expression, and so is `x $\to$f (x)', while `2 + $\times$3' is incorrect and does not represent an object.

It would seem that any correct expression should have---in principle, at least---a unique value, representing some agreed `fully simplified' form of the expression. For instance, it could be argued that the value of 13/91 is 1/7, and that of

$\surd$

2187 is 27

$\surd$

3. The simplified value is more concise and informative, and would enable us---among other things---to recognise when two expressions represent the same object.

This is not so simple. For example, the two expressions

$\surd$ 3 $-$

$\surd$

$\surd$ 5 $-$2

2

6

have the same value, yet there is no compelling reason for choosing one over the other. The same object may be viewed from different angles, and our choice of representation will depend on the context.

The following well-known identity makes an even stronger point:

1 $-$xn

1 $-$x = 1 + x + x2 + \textperiodcentered{} \textperiodcentered{} \textperiodcentered{} + xn$-$1.

The right-hand side---the sum of n monomials---is the `fully simplified' version, while the `unsimplified' left-hand side has only four terms. While simplification reveals thetruenatureof theobject---apolynomial as opposedtoarational function--- it's the unsimplified value which provides information about the sum.

3.4 Expressions 41

Given that agreeing on a unique value of an expression proves difficult, we shift our attention to the type of value (set, number, function, etc.), which assigns the expression to a certain class of objects. This class provides the broadest characterisation of the object in question. Again, we must exercise some judgement. The expressions

$\pi$

sin(x)dx

1 + 1

0

have the same value, but their structure is so different that the coincidence of their values seems secondary. Whereas the expression on the left is unquestionably `a number', or `a positive integer', that on the right is `a definite integral.' On the other hand, there may be circumstances in which the reductionist description of the integral as a number is appropriate, for instance when discussing integrability of functions.

Keeping in mind the difficulties mentioned above, we now turn to the description of mathematical expressions.

\section{3.4.1 Levels of Description}

We develop the idea of successive refinements in the description of an expression, from the general to the particular. The appropriate level of details to be included will vary,dependingonthesituation.Wetreatinparallelverbalandsymbolicdescriptions, as far as it is reasonable to do so.

Let us consider the definition of a set. The coarsest level of description is

\{ . . . \} A set

wheretheobject'stypeisidentifiedbythecurlybrackets.Theuseoftheindeterminate article---`a' set rather than `the' set---reflects our incomplete knowledge.

The next level in specialisation identifies the ambient set:

\{(x, y) $\in$Z2 : . . . \} A set of integer pairs

Now we begin to build the defining properties of our set:

\{(x, y) $\in$Z2 : gcd(x, y) = 1, . . . \} A set of pairs of co-prime integers

The final step completes the definition:

\{(x, y) $\in$Z2 : gcd(x, y) = 1, 2|xy\} The set of pairs of co-prime integers,

with exactly one even component

42 3 Essential Dictionary II

Accordingly, the indefinite article has been replaced by the definite article. Now both words and symbols describe one and the same object, and one should consider the relative merits of the two presentations. The term `exactly' is, strictly speaking, redundant, but it helps the reader realise that x and y cannot both be even. A robotic translation of symbols into words

The set of elements of the cartesian product of the integers with themselves, whose components have greatest common divisor equal to 1, and such that 2 divides the product of the components

lacks the synthesis that comes with understanding.

Expressions may be nested, like boxes within boxes. We begin with the two expressions

2

\textperiodcentered{} \textperiodcentered{} \textperiodcentered{} A square

\textperiodcentered{} \textperiodcentered{} \textperiodcentered{} A sum

We only see the outer structure of these objects. We compose them in two different ways:

2

\textperiodcentered{} \textperiodcentered{} \textperiodcentered{}

The square of a sum

2

\textperiodcentered{} \textperiodcentered{} \textperiodcentered{} A sum of squares

The first term in each expression identifies the object's outer layer. There is still one indefinite article, reflecting a degree of generality. We specialise further:

1

2

$\infty$

The sum of the square of the reciprocal of the

n

natural numbers

n=1

Words or symbols now define a unique object, with three levels of nesting. By contrast, in the nested expression

2

$\infty$

The square of the sum of the elements of

an $\in$Q

an

a rational sequence

n=1

the innermost object---a rational sequence---is still generic.

In these examples words and symbols are interchangeable; in actual writing, some concepts are best expressed with words, others with symbols, while most of them require both. For instance, the symbolic expression

(1 $-$x, 2x + x2, . . . , nxn$-$1 + ($-$x)n, . . .) (3.20)

3.4 Expressions 43

defines an infinite sequence succinctly and unambiguously. Using words, we could begin to describe it as follows:

A sequence An infinite sequence An infinite sequence of binomials

Increasing further the accuracy of the verbal description is pointless, since the symbolic expression (3.20) is clearly superior in delivering exact information. On the other hand, words can place this expression in a context, which is something symbols cannot do. To illustrate this point, we supplement the description given above with additional information which emphasises a particular property.

An infinite sequence of binomials

with integer coefficients with unbounded coefficients with increasing degree whose leading term alternates in sign.

\section{3.4.2 Describing Expressions}

We expand our dictionary with terms describing broad attributes of expressions.

An expression involving numbers, the four arithmetical operations, and raising to an integer or fractional power (extraction of roots), is called an arithmetical expression. The value of an arithmetical expression is a number. A combination of rational numbers and square roots of rational numbers is called a quadratic irrational or a quadratic surd. The following expressions are arithmetical.

1918612 $-$3 \textperiodcentered{} 1107712 = $-$2 An arithmetic identity, with a surprising

cancellation.

$\surd$ 3 + 2

2

$\surd$ 8 $-$3

The ratio of two quadratic surds having 7 distinct radicands.

If indeterminates are present, we speak of an algebraic expression.

6

ab $-$(ab)$-$1

3$\surd$ a2b2 + ab + 1

An algebraic expression with two indeterminates and higher-index roots.

Polynomials and rational functions are algebraic expressions (see Sect. 3.1). They may also be characterised as rational expressions, since they don't involve fractional powers of the indeterminates.

44 3 Essential Dictionary II

1 1 1 x + 1 + x2 + 1 + \textperiodcentered{} \textperiodcentered{} \textperiodcentered{} + xn + 1 The sum of finitely many rational

functions, with increasing degree.

2

2 + 1

x2 + 1

A rational expression, involving repeated compositions of a polynomial function with itself.

+ 1

(x2 + 1)2 + 1

The following mathematical Russian doll

3x + \textperiodcentered{} \textperiodcentered{} \textperiodcentered{} + $\surd$nx n $\in$N

x +

2x +

could be described as

An algebraic expression in one indeterminate, featuring a finite number of nested square roots.

The functions sine, cosine, tangent, secant, etc., are called trigonometric functions (or circular functions). A trigonometric expression is an expression containing trigonometric functions.

8 cos(z)4 $-$8 cos(z)2 + 1 A quartic trigonometric polynomial.

Trigonometric functions belong to the larger class of transcendental functions, which are functions not definable by an algebraic expression (the exponential, the logarithm, etc.). In the expressions (3.16), we have used the terms algebraic, trigonometric and transcendental to describe equations.

The term analytical expression is used in the presence of infinite processes.

3$\surd$ 1 + x = 1 + 1 3x $-$1 9x2 + 5 81x3 + \textperiodcentered{} \textperiodcentered{} \textperiodcentered{} The first few terms of the series expansion of an algebraic function.

n

Napier's constant as the limit of a = e rational sequence.

1 + 1

lim

n

n$\to$$\infty$

$\infty$

(2k)2

An infinite product formula for (2k)2 $-$1 = $\pi$ Archimedes$'$ constant.

2

k=1

An integral expression is an expression involving integrals.

3.4 Expressions 45

x

1 t dt An integral expression for the natural logarithm.

ln(x) =

1

The term combinatorial is appropriate for expressions involving counting functions, such as the factorial function, or the binomial coefficient.

1 (2k)!

A rational combination of exponentials and factorials. 22k (k!)2

n + k

n

A finite sum of binomial coefficients.

k

k=0

In Chap. 4 we deal with the expressions found in logic: the boolean expressions.

Exercise 3.2 For each expression, provide two levels of description: [ $\varepsilon$ ]

(i) a coarse description, which only identifies the object's type (set, function,

equation, etc); (ii) a finer description, which defines the object in question, or characterises its

structure.

1. 33 + 43 + 53 = 63

3$\surd$ 4$\surd$

$\surd$ 2.

3 + $\surd$ 3 + 3

3. 7 5 <

2 < 17

12

4. x3 $-$x $-$1

5. xy > 0

6. (x + y)3 = x3 + 3x2y + 3xy2 + y3

7. y $-$x2 $-$x = 0

8. 1 = 0

9. sin(x $-$y) = sin(x) cos(y) $-$cos(x) sin(y)

10. d2x dt2 $-$2dx

dt $-$3 = 0

11. f (x, y) $\leq$g(x, y)

f (x, y) = 0

12.

g(x, y) = 0

13. (A $\cup$B) $\cap$C = (A $\cap$C) $\cup$(B $\cap$C)

14 + 1 1 24 + 1 34 + \textperiodcentered{} \textperiodcentered{} \textperiodcentered{} = $\pi$4 14.

90 .

Exercise 3.3 Same as in Exercise 3.2.

1. 2Z$\setminus$4Z

2. \{\{$\emptyset$\}\}

3. (Q$\setminus$Z)2

4. R2$\setminus$Q2

5. (R $\times$ Z) $\cup$(Z $\times$ R)

46 3 Essential Dictionary II

6. 1 + 2(1 + 2(1 + 2Z))

7. (\{a\}, \{a\}, \{a\}, . . .)

8. (1, $-$3, 5, . . . , ($-$1)n(2n + 1), . . .)

9. (1, 22k, 32k, . . . , n2k)

10. \{(1, 3), (3, 5), (5, 7), . . .\}

11. ((x1), (x1, x2), (x1, x2, x3), . . .)

12. (\{$\omega$\}, \{\{$\omega$\}\}, \{\{\{$\omega$\}\}\}, . . .).

Exercise 3.4 Same as in Exercise 3.2.

1. f : R $\to$R, x $\to$x + 1

$\infty$

f (x) g(x) dx

2.

0

dx f (x)g(x) = d f (x) d 3.

g(x) + dg(x) dx

f (x) dx

x

d f (y)

4. f (x) =

dy dy

0

F(x, y)dx

5.

1

1

H(x, y)dy

dx

6.

0

0

$\infty$

($-$1)kx2k

7. cos(x) =

(2k)!

k=0

8. $\partial$F(x, y, z) + $\partial$F(x, y, z)

+ $\partial$F(x, y, z) $\partial$y

$\partial$x

$\partial$z

n

$\partial$L ($\theta$1, . . . , $\theta$n)

9.

$\partial$$\theta$k

k=1

xn2

10.

k$\geq$1

$\infty$

xt$-$1e$-$xdx

11.

0

1 $-$x2

.

12. $\pi$ x

n2

n$\geq$1

3.5 Some Advanced Terms 47

\section{3.5 Some Advanced Terms}

We introduce some advanced words and symbols concerning sets, sequences, and equations.

\section{3.5.1 Sets and Sequences}

Let (Ak) be a sequence of sets, which may be finite or infinite. The binary set operations of union and intersection generalise to an arbitrary number of operands:

Ak = A1 $\cup$A2 $\cup$\textperiodcentered{} \textperiodcentered{} \textperiodcentered{}

Ak = A1 $\cap$A2 $\cap$\textperiodcentered{} \textperiodcentered{} \textperiodcentered{} .

k

k

These expressions denote the set of elements belonging to at least one of the sets

Ak, and to every one of the sets Ak, respectively. The integer index conforms to the notational rules for sums and products (Sect. 3.2):

$\infty$

$\Theta$k

(Ak$\setminus$Bk).

k=2

k$\in$2N k$\geq$K

However, infinite unions and intersections may also be controlled by a real (as opposed to integer) index, which adds flexibility to this notation:

$\Omega$$\alpha$ $\Omega$$\alpha$ = \{(x, y) $\in$R2 : y = x2 + $\alpha$x\}.

$\Omega$ =

$\alpha$>0

(Sketch the set $\Omega$.)

A sequence of sets is descending (or nested) if

A1 $\supset$A2 $\supset$A3 $\supset$\textperiodcentered{} \textperiodcentered{} \textperiodcentered{}

\section{and ascending if}

A1 $\subset$A2 $\subset$A3 $\subset$\textperiodcentered{} \textperiodcentered{} \textperiodcentered{} .

Example. The following symbols

\{0\} if |m| = 1 Z if |m| = 1

mZ $\supset$m2Z $\supset$\textperiodcentered{} \textperiodcentered{} \textperiodcentered{} $\supset$mkZ $\supset$\textperiodcentered{} \textperiodcentered{} \textperiodcentered{}

mkZ =

k$\geq$1

tell a story. Say it with words. [ $\varepsilon$ ]

48 3 Essential Dictionary II

The set of multiples of a power of an integer contains the set of multiples of any higher power of the same integer. By considering increasing powers we obtain an infinite nested sequence of sets. Apart from a trivial case, the only element common to all these sets is the origin.

Example. The Farey sequence4 (Fn) is an ascending sequence of sets. Its general term Fn is the set of all reduced fractions in the unit interval whose denominator does not exceed n:

F1 = \{0, 1\}

0, 1 2, 1

F2 =

0, 1 3, 1 2, 2 3, 1

F3 =

0, 1 4, 1 3, 1 2, 2 3, 3 4, 1

F4 =

...

Let us now consider the representation of sets of sequences. Let A be any set. We form the set of all finite sequences of elements of A, with n terms. This set is the cartesian product An of n copies of the set A: the first element a1 of a sequence is chosen from the first copy of A, the second element from the second copy of A, and so on. For instance, the set \{0, 1\}n is the set of all binary sequences with n-digits, while Qn is the set of all n-uples of rational numbers. It follows that the infinite union

Zn

n$\geq$1

represents the set of all finite integer sequences.

The set of all infinite sequences of elements of A has the structure of a cartesian product with infinitely many terms. It would seem natural to denote it by A$\infty$. However, the idiomatic notation AN is more common, due to its greater flexibility. Thus

AZ denotes the set of doubly-infinite sequences of elements of A, and one even finds AZ2 for the set of two-dimensional arrays: (ai, j) $\in$AZ2.

\section{3.5.2 More on Equations}

An equation whose unknown is a function is called a functional equation.

f (x) = f (x + 1) f ( f (x)) = x. (3.21)

4 This sequence, named after John Farey, Sr. (English: 1766--1826), was actually discovered by the Frenchman Charles Haros.

3.5 Some Advanced Terms 49

Here the unknown is f , not x. To make this plain, let 1 be the identity function in the appropriate set of functions, and let $\tau$ be the function that increases its argument by 1: $\tau$(x) = x + 1. Equations (3.21) may now be rewritten without reference to the function's argument, using the composition operator---see (2.19).

f = f $\tau$ f f = 1. (3.22)

Solutions of these equations are periodic functions and involutions, respectively.

The class of functional equations includes the differential equations (which we've seen in Sect.3.3) and the integral equations. The latter are equations in which the unknown function appears under an integral sign, e.g.,

$\infty$

$\surd$$\pi$ f (x) =

e$-$xy f (y)dy.

0

An equation whose unknown is a set is called a set equation:

(X + X) $\cap$X = $\emptyset$. (3.23)

Here the symbol X necessarily represents a set, and hence X + X is an algebraic sum of sets [defined in (2.21)]. This set equation may be defined over any family of sets of numbers, such as P(Z), the power set of the integers. A solution of this equation is called a sum-free set; thus the set of odd integers is sum-free.

It should be clear that equations can be written having any type of mathematical object as unknown.

Exercise 3.5 Same as in Exercise 3.2.

f $-$1

Ak

1.

k

Pj,k

2.

j,k

=

$\Theta$$\alpha$

f ($\Theta$$\alpha$)

3. f

$\alpha$

$\alpha$

$\surd$

2Z)n

4. (Z +

5. \{1, x, x2\}N

6. ZZ

7. X = \{X\}

8. (z + 1) = z(z)

= 1

9. f f \textperiodcentered{} \textperiodcentered{} \textperiodcentered{} f

n

10. [0, 1] + [0, 1] + \textperiodcentered{} \textperiodcentered{} \textperiodcentered{} + [0, 1]

.

n

50 3 Essential Dictionary II

Exercise 3.6 Write an equation whose unknown is (i) a matrix; (ii) a polynomial; (iii) a sequence.

Exercise 3.7 Ask some questions about the Farey sequence.

Exercise 3.8 Let X, Y be sets, and let F, G : X $\to$P(Y). Prove that the set equation F(x) = G (x) is equivalent to the equation

F(x) $\Delta$ G (x) = $\emptyset$

in the sense that the two equations have the same solution set. Thus every set equation may be reduced to the form F(x) = $\emptyset$, which resembles (3.17).

\end{document}